%----------------------------------------------------------------------------------------------------
% START OF CHAPTER
%----------------------------------------------------------------------------------------------------
\chapter{Jan 17}
\paragraph{Topics}
\textit{Nouns and adjectives. Quantitative adjectives. Negation. Helper verbs. Greetings.}

%BIB try this: \cite{Mango01} and see what happens
\section{Grammar}
\subsection{Nouns - gender and articles}

Italian nouns are gendered, masculine or feminine. Nouns ending \iti{-o} are masculine; 
those ending \iti{-a}, feminine. For example, (the) table, \iti{il tavolo}, or (the) house, 
\iti{la casa}. A few nouns end in \iti{-e}, e.g., \iti{il caffè} or \iti{la stazione}.

Articles match the gender of the noun. A notable rule is that nouns that start with \iti{z} 
or the combination \iti{s +} consonant have different articles from the usual \iti{un/il/la-}, 
as illustrated in Table \ref{tab:tArticles}.

\begin{table}[ht]
\centering
    \begin{tabular}{ r l  | r l | l }
         \textbf{masc.} & & \textbf{fem.} &  & Notes  \\
         \hline
        a table    & \iti{un tavolo}  & a house    & \iti{una casa}  &  \\
        the table  & \iti{il tavolo}  & the house  & \iti{la casa} &  \\
        the tables & \iti{i tavoli}   & the houses & \iti{le case} &  \\
	 \hline
        an uncle     &  \iti{uno zio} & an aunt & \iti{una zia} &  Nouns starting \iti{z-}\\
        the uncle    &  \iti{il zio}  & the aunt & \iti{la zia} &  \\
        the uncles   &  \iti{gli zii} & the aunts & \iti{le zie} &  \\
         \hline
        a receipt    & \iti{uno scontrino} & a station & \iti{una stazione} & Nouns starting \iti{s+consonant}\\
        the receipt  & \iti{lo scontrino} & the station & \iti{la stazione} \\
        the receipts & \iti{gli scontrini} & the stations & \iti{le stazione}
\end{tabular}
    \caption{Articles}
    \label{tab:tArticles}
\end{table}
\subsection{Adjectives}
    Adjectives have a base form ending with \iti{-o} or \iti{-e}. The former change their 
	ending to \iti{-a} when used with feminine singular nouns. For example \iti{rosso}, 
	red, is \iti{rosso} for masculine and \iti{rossa} for feminine singular nouns. 
	On the other hand, an \iti{-e} adjective like \iti{grande} is just \iti{grande} in 
	the singular for both genders of noun. Whatever the adjective, there must be consistency 
	between the noun and adjective in terms of gender (masculine or feminine) and number 
	(singular or plural). This includes the speaker when the adjective applies to them 
	(i.e., the speaker is the noun).

\paragraph{Examples}
\begin{itemize}
    \item The gelato is good. \\
    \iti{Il gelato è buon\uline{o}.}
    \item The pizza is good. \\
    \iti{La pizza è buon\uline{a}.}
    \item Claudio is tired. \\
    \iti{Claudio è stanc\uline{o}.}
    \item Francesca is glad. \\
    \iti{Francesca è content\uline{a}.}
    \item The coffee is good. \\
    \iti{Il caffè è buono.}
    \item The pizza is not good. \\
    \iti{La pizza non è buona.}    
\end{itemize}
Most adjectives go after the noun, but adjectives of quantity (\iti{molto}), possession 
(\iti{mio, tuo, sua, etc.}) go before. Some adjectives that normally go after can go before 
to stress a point or to subtly alter the sense of the phrase. For example, compare 
``My car is new'' with ``My \textit{new} car''. The adjective \iti{grande} is a case in point: 
when used after a noun, it describes `big'-ness. Used before a noun, it describes `great'-ness, 
as in, ``London is a big city'' vs. ``London is a great city''.

\paragraph{Examples}
\begin{itemize}
    \item The red house. \\
    \iti{La casa rossa.}
    \item The new car. \\
    \iti{La macchina nuova.}
    \item The \iti{new} car. \\
    \iti{La nuova macchina.}
    \item My wife. \\
    \iti{La mia moglie.}
    \item His daughter. \\
    \iti{La sua figlia.}
    \item Rome is a big city. \\
    \iti{Roma è una città grande.}
    \item Rome is a great city. (As in, historic, magnificent, special)\\
    \iti{Roma è una grande città.}
    \item Un tavolo caro. \\
    \iti{An expensive table.}
    \item Umberto is my dear uncle. (Affection)\\
    \iti{Umberto è mio caro zio.}
\end{itemize}
\begin{mdframed}
    Italian seems to have a similar `BANGS' rule to French: adjectives go after the noun 
	unless they talk about \textbf{B}eauty, \textbf{A}ge, \textbf{N}umbers, 
	\textbf{G}oodness, or \textbf{S}ize. See the Appendix.
\end{mdframed}
\subsection{Quantitative adjectives - \iti{molto}}
The word \itb{molto} can act as an adjective, preceding a noun (``I have \textit{a lot of} 
	patience'', ``I don't do \textit{much} sport''). In these cases it \uline{changes 
	to match the noun}.

However, \iti{molto} can also act an adverb, (``Ravenna \textit{is very} beautiful''), 
	and an adverb-like modifier to a verb (``I like this car \textit{a lot}''). 
	In these cases, it does \uline{not} need to change.
\paragraph{Examples}
\begin{itemize}
    \item I talk a lot. \\
    \iti{Io parlo molto.} \\
    No change to \iti{molto}. It is being used as an adverb to \iti{parlo}.

    \item I don't talk much. \\
    \iti{Io non parlo molto.} \\
    No change to \iti{molto}. It is being used as an adverb to \iti{parlo}.

    \item A lot of pasta, please. \\
    \iti{Molt\uline{a} pasta, per favore.} \\
    Adjective for \iti{pasta}, so \iti{molto} must change to \iti{molta} to match gender.
    
    \item (Claudio says \ldots) I am very tired. \\
    \iti{Io sono molt\uline{o} stanco.} \\
    No change to \iti{molto}. It is an adverb to \iti{sono}.

    \item (Francesca says \ldots) I am very tired. \\
    \iti{Io sono molt\uline{o} stanca.} \\
    No change to \iti{molto}. It is an adverb to \iti{sono}. (Notice that \iti{stanca} does 
		take the \iti{-a} ending, as it is an adjective that applies to Francesca.)

    \item I have a lot of dogs. \\
    \iti{Io ho molt\uline{i} cani.} \\
    Adjective for \iti{cani}, so \iti{molti} is required to match the plurality of \iti{cane}.
    
    \item Italy is very beautiful. \\
    \iti{L'Italia è molto bella.} \\
    Adverb for \iti{è}, so no change to \iti{molto} needed.
    
    \item The pizza is very good. \\
    \iti{La pizza è molto buona.} \\
    Adverb for \iti{è}, so no change to \iti{molto} needed.

\end{itemize}

\subsection{The simple negation - \iti{non}}
Statements can be negated by putting \iti{non} between the subject and the verb. 
\paragraph{Examples}
\begin{itemize}
    \item Francesca is not tired. \\
    \iti{Francesca non è stanca.}
    \item The pizza is not good. \\
    \iti{La pizza non è buona.}
    \item I am not tired, because I am not at work today. \\
    \iti{Io non sono stanco, perché io non sono a lavoro oggi.}
\end{itemize}

\section{Vocabulary}
\subsection{Verbs}
\subsection{The basics of \iti{essere} and \iti{avere}}
Table \ref{tab:tVerbsPresentEssereAvere} has the basics.
\begin{table}[ht]
\centering
    \begin{tabular}{ r l | r l }
         \itb{essere} & \iti{to be} & \itb{avere} & \iti{to have} \\
         \hline
         \itb{io}      & \iti{sono} & \itb{io}      & \iti{ho}  \\
         \itb{tu}      & \iti{sei}  & \itb{tu}      & \iti{hai} \\
         \itb{lui/lei} & \iti{è}    & \itb{lui/lei} & \iti{ha}  \\
    \end{tabular}
    \caption{\iti{essere / avere}, the basics, present tense}
    \label{tab:tVerbsPresentEssereAvere}
\end{table}

\subsection{Words}
See Table \ref{tab:tVocabBasics}. (Note: in the table, and elsewhere in this document, 
gendered terms such as adjectives are presented in their masculine form by default.)

\begin{table}[ht] %[ht] % place table roughly [h]ere or else at the [t]op of page
\centering
    \begin{tabular}{ r l  | r l  | r l  }
          table & \iti{tavolo} &  kitchen table & \iti{tavola} &  wine & \iti{vino} \\
          pizza & \iti{pizza} & house/home & \iti{casa}   \\
          \\
          good (thing) & \iti{buono} & good (performance) & \iti{bravo} &  happy & \iti{felice} \\
          glad & \iti{contento} & a lot of & \iti{molto}  & tired  & \iti{stanco}
    \end{tabular}
    \caption{Some basic nouns and adjectives}
    \label{tab:tVocabBasics}
\end{table}

\subsection{Phrases - greetings}
Greetings reflect the level of politeness, custom, and formality. See Table \ref{tab:tGreetings}:

\begin{table}[ht]
\centering
    \begin{tabular}{ r l  | r l  | l }
         \textbf{informal} &                  & \textbf{formal}  &                      & \textbf{Notes}  \\
         \hline
         Hi                & \iti{Ciao}       & Hello            & \iti{Salve}          &                 \\
         Good morning      & \iti{Buongiorno} & Good morning     & \iti{Buongiorno}     &                 \\
         \ldots night      & \iti{Buonanotte} & \ldots night     & \iti{Buonanotte}     &                 \\
                           &                  & \ldots afternoon & \iti{Buonpomeriggio} &                 \\
                           &                  & \ldots evening   & \iti{Buonasera}      &                 \\
         Bye               & \iti{Ciao, ciao} & Goodbye          & \iti{Arrivederci}    & To one or more people \\
                           &                  &                  &  \iti{Arrivederla}   & More formal, to one person \\
                           &                  & Have a good day  & \iti{Buona giornata} &                 \\
                           &                  & \ldots night     & \iti{Buona serata}   & 
    \end{tabular}
    \caption{Greetings}
    \label{tab:tGreetings}
\end{table}
%----------------------------------------------------------------------------------------------------
% END OF CHAPTER
%----------------------------------------------------------------------------------------------------

